\documentclass[a4paper]{article}
\usepackage[pdftex]{graphicx}
\usepackage[utf8]{inputenc}
\usepackage{enumerate}
\usepackage{amssymb}
\usepackage{icomma}
\usepackage{href-ul}
\hypersetup{
	colorlinks=true,
	linkcolor=blue,
	urlcolor=blue}
\usepackage{siunitx}
\sisetup{locale=DE}
\usepackage{geometry}
\geometry{a4paper, top=15mm, left=15mm, right=15mm, bottom=15mm,
	headsep=10mm, footskip=12mm}

\begin{document}
	{\bf density, volume and mass}
	
	\begin{enumerate}[1.]
		{\bf \item Fill in the table.} Note: the density is constant for every substance. The volume and the corresponding mass can be selected as desired. The state is either solid or liquid or gaseous.
		
		\renewcommand{\arraystretch}{2}
		\begin{tabular}{|p{1.7 cm}|p{1.7 cm}|p{1.7 cm}|p{1.7 cm}|p{1.7 cm}|p{1.7 cm}|p{1.7 cm}|p {1.7cm}|}
			\hline
			Cloth & Wood & Lead & Water & Air & Edible Oil & Gold & Aluminum \\
			\hline
			Condition &  &  &  &  &  &  &  \\
			\hline
			Example & Building Block & Fishing Lead & Bottle & Balloon & Bottle & Ingots & Sharpener\\
			\hline
			Mass & $\SI{35}{\gram}$ & $\SI{5,0}{\gram}$ & & $\SI{5,2}{\gram}$ & $\SI{0.63 }{\kilogram}$ & &$\SI{4,3}{\gram}$ \\
			\hline
			Volume &$\SI{50}{\centi \meter^3}$ & &$\SI{1.5}{\deci \meter^3}$ & & $\SI{0.7}{\deci \meter^ 3 }$ &$\SI{0.64}{\deci\meter^3}$ & $\SI{1.6}{\centi\meter^3}$ \\
			\hline
			Density & &$\SI{11,3}{\gram\over{\centi\meter^3}}$ &$\SI{1,0}{\gram\over{\centi\meter^3}}$ &$\SI{1.3}{\gram\over{\deci \meter^3}}$ & &$\SI{19,3}{\gram\over{\centi \meter^3}}$ & \\
			\hline
		\end{tabular}
		\vspace{0.5 cm}
		
		\begin{tabular}{|p{1.4 cm}|p{1.4 cm}|p{1.5 cm}|p{1.8 cm}|p{1.7 cm}|p{2.0 cm}|p{2.1 cm}|p {1.7cm}|}
			\hline
			Fabric & Cork & Iron & Styrofoam & Hydrogen & Iridium & Mercury & Polyethylene \\
			\hline
			State & & & & & fixed & & \\
			\hline
			Example & cork & key & insulation board & balloon & original kilogram & thermometer & dice\\
			\hline
			Mass &$\SI{1.9}{\gram}$ & & $\SI{150}{\gram}$ & $\SI{0.41}{\gram}$ & & $\SI{0, 67}{\gram}$ & $\SI{3,1}{\gram}$\\
			\hline
			Volume & & $\SI{1,1}{\centi\meter^3}$ & $\SI{5,0}{\deci\meter^3}$ & & $\SI{45}{\centi\meter^3}$ &$\SI{0.05}{\centi\meter^3}$ & \\
			\hline
			Density &$\SI{0.3}{\gram\over{\centi \meter^3}}$ &$\SI{7.7}{\gram\over{\centi \meter^3}}$ & &$\SI{90}{\gram\over{\meter^3}}$ &$\SI{22,4}{\gram\over{\centi\meter^3}}$ & &$\SI{ 0.95}{\gram\over{\centi\meter^3}}$ \\
			\hline
		\end{tabular}
		
		{\bf \item Which is the specifically lightest material?} Which is the one with the highest density?
		
		{\bf \item Why does a hydrogen balloon fly up?} Its weight is not negative...
		
		{\bf \item task}
		
		\begin{enumerate}[a)]
			\item Which of the solid substances mentioned above would be able to float in which liquid?
			
			\renewcommand{\arraystretch}{4}
			\begin{tabular}{|p{3 cm}|p{12.5 cm}|}
				\hline
				Floats in & Substance\\
				\hline
				cooking oil &\\
				\hline
				Water	& \\
				\hline
				Mercury & \\
				\hline
			\end{tabular}
			
			\item Which substance floats in water but sinks in cooking oil?
			\item Which substances are lost in mercury?
		\end{enumerate}
	\end{enumerate}
	\fbox{
		\begin{minipage}{0.5\textwidth}
			For the solution, please \href{https://www.okuyakl.de/phys/p7dihL073/le073.pdf}{click here} or scan the QR code.\\
			Additional worksheets are available at
			
			\href{http://www.okuyakl.com}{www.okuyakl.com}
		\end{minipage}
		\hfill
		\begin{minipage}{0.4\textwidth}
			\includegraphics[width=1.5 cm]{../../viecher/zwe03}
			\includegraphics[width=3 cm]{qre073}
			\includegraphics[width=2 cm]{../../viecher/afanticon1}
			
	\end{minipage}}
	
\end{document}